\documentclass[12pt]{article}
\usepackage[UTF8]{inputenc}
\usepackage{xeCJK}
\setCJKmainfont{Sarabun}
\usepackage{enumitem}
\usepackage{geometry}
\geometry{a4paper, margin=2.5cm}
\usepackage{setspace}
\onehalfspacing

\begin{document}

\begin{center}
{\Large\textbf{ฟิสิกส์อะตอม}}\\[6pt]
{\normalsize แบบฝึกหัดเพิ่มเติม}\\[6pt]
{\normalsize วิชาฟิสิกส์ ม.ปลาย บทที่ 25}\\[4pt]
{\footnotesize ทั้งหมด 30 ข้อ}
\end{center}
\hrule
\bigskip

\section*{แบบฝึกหัดเพิ่มเติม}
\begin{enumerate}[label=\arabic*., itemsep=0.5\baselineskip, topsep=0.5\baselineskip, leftmargin=*]
\item ใครค้นพบอิเล็กตรอนและเสนอแบบจำลองลูกกลมโป้ยปุ้น?
\begin{enumerate}[label=\alph*)., itemsep=0.3\baselineskip, topsep=0.3\baselineskip]
\item ก\ 
\begin{minipage}[t]{0.9\linewidth}
\raggedright Dalton
\end{minipage}
\item ข\ 
\begin{minipage}[t]{0.9\linewidth}
\raggedright Thomson
\end{minipage}
\item ค\ 
\begin{minipage}[t]{0.9\linewidth}
\raggedright Rutherford
\end{minipage}
\item ง\ 
\begin{minipage}[t]{0.9\linewidth}
\raggedright Bohr
\end{minipage}
\end{enumerate}
\vspace{-0.3\baselineskip}
\begin{small}
\textit{คำตอบ: ข}
\end{small}

\item ข้อใดเกี่ยวข้องกับการทดลองแผ่นทองคำของรัทเทอร์ฟอร์ดโดยตรง?
\begin{enumerate}[label=\alph*)., itemsep=0.3\baselineskip, topsep=0.3\baselineskip]
\item ก\ 
\begin{minipage}[t]{0.9\linewidth}
\raggedright  cathode ray tube
\end{minipage}
\item ข\ 
\begin{minipage}[t]{0.9\linewidth}
\raggedright อนุภาคแอลฟา
\end{minipage}
\item ค\ 
\begin{minipage}[t]{0.9\linewidth}
\raggedright สเปกตรัมเส้น
\end{minipage}
\item ง\ 
\begin{minipage}[t]{0.9\linewidth}
\raggedright กฎของนิวตัน
\end{minipage}
\end{enumerate}
\vspace{-0.3\baselineskip}
\begin{small}
\textit{คำตอบ: ข}
\end{small}

\item สถานะพื้น (ground state) ของอิเล็กตรอนในอะตอมไฮโดรเจนมีเลขควอนตัมหลัก n เท่ากับเท่าใด?
\begin{enumerate}[label=\alph*)., itemsep=0.3\baselineskip, topsep=0.3\baselineskip]
\item ก\ 
\begin{minipage}[t]{0.9\linewidth}
\raggedright 0
\end{minipage}
\item ข\ 
\begin{minipage}[t]{0.9\linewidth}
\raggedright 1
\end{minipage}
\item ค\ 
\begin{minipage}[t]{0.9\linewidth}
\raggedright 2
\end{minipage}
\item ง\ 
\begin{minipage}[t]{0.9\linewidth}
\raggedright ∞
\end{minipage}
\end{enumerate}
\vspace{-0.3\baselineskip}
\begin{small}
\textit{คำตอบ: ข}
\end{small}

\item เมื่ออิเล็กตรอนกระโดดจากระดับ n=4 ลง n=2 จะปล่อยโฟตอนที่มีพลังงานเท่าใด? (E\_n = -13.6/n² eV)
\begin{enumerate}[label=\alph*)., itemsep=0.3\baselineskip, topsep=0.3\baselineskip]
\item ก\ 
\begin{minipage}[t]{0.9\linewidth}
\raggedright 0.66 eV
\end{minipage}
\item ข\ 
\begin{minipage}[t]{0.9\linewidth}
\raggedright 2.55 eV
\end{minipage}
\item ค\ 
\begin{minipage}[t]{0.9\linewidth}
\raggedright 3.4 eV
\end{minipage}
\item ง\ 
\begin{minipage}[t]{0.9\linewidth}
\raggedright 12.75 eV
\end{minipage}
\end{enumerate}
\vspace{-0.3\baselineskip}
\begin{small}
\textit{คำตอบ: ข}
\end{small}

\item สเปกตรัมเส้นของธาตุใช้ระบุธาตุได้เพราะเหตุผลใด?
\begin{enumerate}[label=\alph*)., itemsep=0.3\baselineskip, topsep=0.3\baselineskip]
\item ก\ 
\begin{minipage}[t]{0.9\linewidth}
\raggedright ทุกธาตุมีสเปกตรัมเหมือนกัน
\end{minipage}
\item ข\ 
\begin{minipage}[t]{0.9\linewidth}
\raggedright สเปกตรัมเฉพาะของแต่ละธาตุเพราะระดับพลังงานเฉพาะ
\end{minipage}
\item ค\ 
\begin{minipage}[t]{0.9\linewidth}
\raggedright ธาตุเบามีสเปกตรัมมากกว่าธาตุหนัก
\end{minipage}
\item ง\ 
\begin{minipage}[t]{0.9\linewidth}
\raggedright สเปกตรัมเส้นเปลี่ยนตามอุณหภูมิ
\end{minipage}
\end{enumerate}
\vspace{-0.3\baselineskip}
\begin{small}
\textit{คำตอบ: ข}
\end{small}

\item แบบจำลองอะตอมแบบใดที่บอกว่าอะตอมเป็นทรงกลมแข็งไม่สามารถแบ่งแยกได้?
\begin{enumerate}[label=\alph*)., itemsep=0.3\baselineskip, topsep=0.3\baselineskip]
\item ก\ 
\begin{minipage}[t]{0.9\linewidth}
\raggedright แบบจำลองของ Thomson
\end{minipage}
\item ข\ 
\begin{minipage}[t]{0.9\linewidth}
\raggedright แบบจำลองของ Dalton
\end{minipage}
\item ค\ 
\begin{minipage}[t]{0.9\linewidth}
\raggedright แบบจำลองของ Rutherford
\end{minipage}
\item ง\ 
\begin{minipage}[t]{0.9\linewidth}
\raggedright แบบจำลองของ Bohr
\end{minipage}
\end{enumerate}
\vspace{-0.3\baselineskip}
\begin{small}
\textit{คำตอบ: ข}
\end{small}

\item ป้ายหลอดไฟนีออนที่เปล่งแสงสีแดง-ส้ม เกิดจากหลักการใด?
\begin{enumerate}[label=\alph*)., itemsep=0.3\baselineskip, topsep=0.3\baselineskip]
\item ก\ 
\begin{minipage}[t]{0.9\linewidth}
\raggedright การเปล่งแสงของอะตอมก๊าซเมื่อถูกกระตุ้นด้วยไฟฟ้า
\end{minipage}
\item ข\ 
\begin{minipage}[t]{0.9\linewidth}
\raggedright การเผาไหม้ของก๊าซนีออน
\end{minipage}
\item ค\ 
\begin{minipage}[t]{0.9\linewidth}
\raggedright การสะท้อนของแสงขาว
\end{minipage}
\item ง\ 
\begin{minipage}[t]{0.9\linewidth}
\raggedright การเรืองแสงของของแข็ง
\end{minipage}
\end{enumerate}
\vspace{-0.3\baselineskip}
\begin{small}
\textit{คำตอบ: ก}
\end{small}

\item รัศมีวงโคจรของอิเล็กตรอนในแบบจำลองโบร์แปรผันกับ n² หมายความว่าอย่างไร?
\begin{enumerate}[label=\alph*)., itemsep=0.3\baselineskip, topsep=0.3\baselineskip]
\item ก\ 
\begin{minipage}[t]{0.9\linewidth}
\raggedright รัศมีเพิ่มขึ้นตาม n
\end{minipage}
\item ข\ 
\begin{minipage}[t]{0.9\linewidth}
\raggedright รัศมีเพิ่มขึ้นเป็น n² เท่า
\end{minipage}
\item ค\ 
\begin{minipage}[t]{0.9\linewidth}
\raggedright รัศมีลดลงเป็น n² เท่า
\end{minipage}
\item ง\ 
\begin{minipage}[t]{0.9\linewidth}
\raggedright รัศมีคงที่ทุกระดับ
\end{minipage}
\end{enumerate}
\vspace{-0.3\baselineskip}
\begin{small}
\textit{คำตอบ: ข}
\end{small}

\item ชุด Paschen ในสเปกตรัมของไฮโดรเจน ให้รังสีในช่วงใดของสเปกตรัมแม่เหล็กไฟฟ้า?
\begin{enumerate}[label=\alph*)., itemsep=0.3\baselineskip, topsep=0.3\baselineskip]
\item ก\ 
\begin{minipage}[t]{0.9\linewidth}
\raggedright UV
\end{minipage}
\item ข\ 
\begin{minipage}[t]{0.9\linewidth}
\raggedright Visible
\end{minipage}
\item ค\ 
\begin{minipage}[t]{0.9\linewidth}
\raggedright Infrared
\end{minipage}
\item ง\ 
\begin{minipage}[t]{0.9\linewidth}
\raggedright Microwave
\end{minipage}
\end{enumerate}
\vspace{-0.3\baselineskip}
\begin{small}
\textit{คำตอบ: ค}
\end{small}

\item พลังงานของโฟตอนที่ปล่อยเมื่ออิเล็กตรอนกระโดดจาก n=5 ลง n=3 มีค่าเท่าใด?
\begin{enumerate}[label=\alph*)., itemsep=0.3\baselineskip, topsep=0.3\baselineskip]
\item ก\ 
\begin{minipage}[t]{0.9\linewidth}
\raggedright 0.31 eV
\end{minipage}
\item ข\ 
\begin{minipage}[t]{0.9\linewidth}
\raggedright 0.97 eV
\end{minipage}
\item ค\ 
\begin{minipage}[t]{0.9\linewidth}
\raggedright 1.51 eV
\end{minipage}
\item ง\ 
\begin{minipage}[t]{0.9\linewidth}
\raggedright 2.86 eV
\end{minipage}
\end{enumerate}
\vspace{-0.3\baselineskip}
\begin{small}
\textit{คำตอบ: ข}
\end{small}

\item แบบจำลองโบร์เสนอเงื่อนไขการโคจรของอิเล็กตรอนว่าอิเล็กตรอนต้องมีขนาดของโมเมนตัมเชิงมุมเป็นทวีคูณของ ħ หมายความว่าอย่างไร?
\begin{enumerate}[label=\alph*)., itemsep=0.3\baselineskip, topsep=0.3\baselineskip]
\item ก\ 
\begin{minipage}[t]{0.9\linewidth}
\raggedright ระดับพลังงานเป็นค่าต่อเนื่อง
\end{minipage}
\item ข\ 
\begin{minipage}[t]{0.9\linewidth}
\raggedright ระดับพลังงานเป็นค่าไม่ต่อเนื่อง (quantized)
\end{minipage}
\item ค\ 
\begin{minipage}[t]{0.9\linewidth}
\raggedright อิเล็กตรอนเคลื่อนที่เป็นเส้นตรง
\end{minipage}
\item ง\ 
\begin{minipage}[t]{0.9\linewidth}
\raggedright อิเล็กตรอนตกลงนิวเคลียร์
\end{minipage}
\end{enumerate}
\vspace{-0.3\baselineskip}
\begin{small}
\textit{คำตอบ: ข}
\end{small}

\item สารเรืองแสง (fluorescent paint) ในเสื้อผ้าสีสันสดใสทำงานอย่างไร?
\begin{enumerate}[label=\alph*)., itemsep=0.3\baselineskip, topsep=0.3\baselineskip]
\item ก\ 
\begin{minipage}[t]{0.9\linewidth}
\raggedright ดูดกลืนแสง UV ในแสงแดดแล้วเปล่งแสงสีที่ตามองเห็นได้
\end{minipage}
\item ข\ 
\begin{minipage}[t]{0.9\linewidth}
\raggedright ดูดกลืนแสงทุกสีแล้วเปลี่ยนเป็นความร้อน
\end{minipage}
\item ค\ 
\begin{minipage}[t]{0.9\linewidth}
\raggedright สะท้อนแสงสีต่างๆ โดยไม่ดูดกลืน
\end{minipage}
\item ง\ 
\begin{minipage}[t]{0.9\linewidth}
\raggedright เปล่งแสงเองตลอดเวลา
\end{minipage}
\end{enumerate}
\vspace{-0.3\baselineskip}
\begin{small}
\textit{คำตอบ: ก}
\end{small}

\item อะตอมไฮโดรเจนอยู่ที่สถานะกระตุ้น (excited state) ที่ n=3 สามารถปล่อยโฟตอนได้กี่วิธี?
\begin{enumerate}[label=\alph*)., itemsep=0.3\baselineskip, topsep=0.3\baselineskip]
\item ก\ 
\begin{minipage}[t]{0.9\linewidth}
\raggedright 1 วิธี
\end{minipage}
\item ข\ 
\begin{minipage}[t]{0.9\linewidth}
\raggedright 2 วิธี
\end{minipage}
\item ค\ 
\begin{minipage}[t]{0.9\linewidth}
\raggedright 3 วิธี
\end{minipage}
\item ง\ 
\begin{minipage}[t]{0.9\linewidth}
\raggedright 4 วิธี
\end{minipage}
\end{enumerate}
\vspace{-0.3\baselineskip}
\begin{small}
\textit{คำตอบ: ค}
\end{small}

\item ค่าคงที่ Rydberg สำหรับไฮโดรเจน (R\_H) มีค่าเท่าใดโดยประมาณ?
\begin{enumerate}[label=\alph*)., itemsep=0.3\baselineskip, topsep=0.3\baselineskip]
\item ก\ 
\begin{minipage}[t]{0.9\linewidth}
\raggedright 1.097 × 10⁵ m⁻¹
\end{minipage}
\item ข\ 
\begin{minipage}[t]{0.9\linewidth}
\raggedright 1.097 × 10⁶ m⁻¹
\end{minipage}
\item ค\ 
\begin{minipage}[t]{0.9\linewidth}
\raggedright 1.097 × 10⁷ m⁻¹
\end{minipage}
\item ง\ 
\begin{minipage}[t]{0.9\linewidth}
\raggedright 1.097 × 10⁸ m⁻¹
\end{minipage}
\end{enumerate}
\vspace{-0.3\baselineskip}
\begin{small}
\textit{คำตอบ: ค}
\end{small}

\item ทำไมแบบจำลองโบร์จึงต้องอาศัยแบบจำลองวงโคจรของอิเล็กตรอน?
\begin{enumerate}[label=\alph*)., itemsep=0.3\baselineskip, topsep=0.3\baselineskip]
\item ก\ 
\begin{minipage}[t]{0.9\linewidth}
\raggedright เพราะอิเล็กตรอนเคลื่อนที่เป็นวงโคจรจริง
\end{minipage}
\item ข\ 
\begin{minipage}[t]{0.9\linewidth}
\raggedright เพราะต้องการอธิบายการแผ่รังสีของอิเล็กตรอนที่โคจร
\end{minipage}
\item ค\ 
\begin{minipage}[t]{0.9\linewidth}
\raggedright เพราะโบร์ต้องการให้สอดคล้องกับฟิสิกส์คลาสสิก
\end{minipage}
\item ง\ 
\begin{minipage}[t]{0.9\linewidth}
\raggedright เพราะอิเล็กตรอนไม่มีโมเมนตัมเชิงมุม
\end{minipage}
\end{enumerate}
\vspace{-0.3\baselineskip}
\begin{small}
\textit{คำตอบ: ค}
\end{small}

\item เส้นสเปกตรัมที่ 4 ของชุด Balmer เกิดจากการกระโดดใด?
\begin{enumerate}[label=\alph*)., itemsep=0.3\baselineskip, topsep=0.3\baselineskip]
\item ก\ 
\begin{minipage}[t]{0.9\linewidth}
\raggedright n=5 → n=2
\end{minipage}
\item ข\ 
\begin{minipage}[t]{0.9\linewidth}
\raggedright n=6 → n=2
\end{minipage}
\item ค\ 
\begin{minipage}[t]{0.9\linewidth}
\raggedright n=7 → n=2
\end{minipage}
\item ง\ 
\begin{minipage}[t]{0.9\linewidth}
\raggedright n=8 → n=2
\end{minipage}
\end{enumerate}
\vspace{-0.3\baselineskip}
\begin{small}
\textit{คำตอบ: ข}
\end{small}

\item อิเล็กตรอนที่ระดับ n มีพลังงานต่ำลงเมื่อ n มีค่าอย่างไร?
\begin{enumerate}[label=\alph*)., itemsep=0.3\baselineskip, topsep=0.3\baselineskip]
\item ก\ 
\begin{minipage}[t]{0.9\linewidth}
\raggedright n ยิ่งน้อย พลังงานยิ่งต่ำ
\end{minipage}
\item ข\ 
\begin{minipage}[t]{0.9\linewidth}
\raggedright n ยิ่งมาก พลังงานยิ่งต่ำ
\end{minipage}
\item ค\ 
\begin{minipage}[t]{0.9\linewidth}
\raggedright n ไม่มีผลต่อพลังงาน
\end{minipage}
\item ง\ 
\begin{minipage}[t]{0.9\linewidth}
\raggedright n ยิ่งมาก พลังงานยิ่งสูง
\end{minipage}
\end{enumerate}
\vspace{-0.3\baselineskip}
\begin{small}
\textit{คำตอบ: ก}
\end{small}

\item จอโทรทัศน์แบบ CRT (Cathode Ray Tube) มีส่วนประกอบที่ทำให้เกิดภาพสีได้อย่างไร?
\begin{enumerate}[label=\alph*)., itemsep=0.3\baselineskip, topsep=0.3\baselineskip]
\item ก\ 
\begin{minipage}[t]{0.9\linewidth}
\raggedright ใช้หลอดฟลูออเรสเซนต์
\end{minipage}
\item ข\ 
\begin{minipage}[t]{0.9\linewidth}
\raggedright ใช้สารเรืองแสง RGB ที่ถูกกระตุ้นด้วยอิเล็กตรอน
\end{minipage}
\item ค\ 
\begin{minipage}[t]{0.9\linewidth}
\raggedright ใช้แสง LED
\end{minipage}
\item ง\ 
\begin{minipage}[t]{0.9\linewidth}
\raggedright ใช้กระจกสะท้อน
\end{minipage}
\end{enumerate}
\vspace{-0.3\baselineskip}
\begin{small}
\textit{คำตอบ: ข}
\end{small}

\item ข้อจำกัดสำคัญของแบบจำลองโบร์คืออะไร?
\begin{enumerate}[label=\alph*)., itemsep=0.3\baselineskip, topsep=0.3\baselineskip]
\item ก\ 
\begin{minipage}[t]{0.9\linewidth}
\raggedright อธิบายไม่ได้ว่าอิเล็กตรอนเคลื่อนที่อย่างไร
\end{minipage}
\item ข\ 
\begin{minipage}[t]{0.9\linewidth}
\raggedright อธิบายสเปกตรัมของอะตอมที่มีหลายอิเล็กตรอนไม่ได้
\end{minipage}
\item ค\ 
\begin{minipage}[t]{0.9\linewidth}
\raggedright ใช้ได้กับทุกธาตุ
\end{minipage}
\item ง\ 
\begin{minipage}[t]{0.9\linewidth}
\raggedright ขัดกับกฎการเคลื่อนที่ของนิวตัน
\end{minipage}
\end{enumerate}
\vspace{-0.3\baselineskip}
\begin{small}
\textit{คำตอบ: ข}
\end{small}

\item ไอออน He⁺ (อะตอมฮีเลียมที่สูญเสียอิเล็กตรอน 1 ตัว) มีระดับพลังงานเป็นอย่างไรเมื่อเทียบกับอะตอมไฮโดรเจน?
\begin{enumerate}[label=\alph*)., itemsep=0.3\baselineskip, topsep=0.3\baselineskip]
\item ก\ 
\begin{minipage}[t]{0.9\linewidth}
\raggedright เหมือนกันทุกประการ
\end{minipage}
\item ข\ 
\begin{minipage}[t]{0.9\linewidth}
\raggedright คล้ายกันแต่คูณด้วย 4
\end{minipage}
\item ค\ 
\begin{minipage}[t]{0.9\linewidth}
\raggedright คล้ายกันแต่คูณด้วย 2
\end{minipage}
\item ง\ 
\begin{minipage}[t]{0.9\linewidth}
\raggedright ไม่มีระดับพลังงาน
\end{minipage}
\end{enumerate}
\vspace{-0.3\baselineskip}
\begin{small}
\textit{คำตอบ: ข}
\end{small}

\item เส้นขีด จำกัด (series limit) ของชุด Balmer คือความยาวคลื่นเมื่อ n₂ → ∞ มีค่าเท่าใด?
\begin{enumerate}[label=\alph*)., itemsep=0.3\baselineskip, topsep=0.3\baselineskip]
\item ก\ 
\begin{minipage}[t]{0.9\linewidth}
\raggedright 364.5 nm
\end{minipage}
\item ข\ 
\begin{minipage}[t]{0.9\linewidth}
\raggedright 400 nm
\end{minipage}
\item ค\ 
\begin{minipage}[t]{0.9\linewidth}
\raggedright 656.3 nm
\end{minipage}
\item ง\ 
\begin{minipage}[t]{0.9\linewidth}
\raggedright 820 nm
\end{minipage}
\end{enumerate}
\vspace{-0.3\baselineskip}
\begin{small}
\textit{คำตอบ: ก}
\end{small}

\item เมื่ออิเล็กตรอนในไฮโดรเจนดูดกลืนโฟตอนพลังงาน 12.09 eV จากสถานะพื้น มันจะอยู่ที่ระดับใด?
\begin{enumerate}[label=\alph*)., itemsep=0.3\baselineskip, topsep=0.3\baselineskip]
\item ก\ 
\begin{minipage}[t]{0.9\linewidth}
\raggedright n=2
\end{minipage}
\item ข\ 
\begin{minipage}[t]{0.9\linewidth}
\raggedright n=3
\end{minipage}
\item ค\ 
\begin{minipage}[t]{0.9\linewidth}
\raggedright n=4
\end{minipage}
\item ง\ 
\begin{minipage}[t]{0.9\linewidth}
\raggedright n=5
\end{minipage}
\end{enumerate}
\vspace{-0.3\baselineskip}
\begin{small}
\textit{คำตอบ: ข}
\end{small}

\item เพชรเรืองแสงสีฟ้าภายใต้แสง UV เพราะเหตุใด?
\begin{enumerate}[label=\alph*)., itemsep=0.3\baselineskip, topsep=0.3\baselineskip]
\item ก\ 
\begin{minipage}[t]{0.9\linewidth}
\raggedright เพชรมีสีฟ้าโดยธรรมชาติ
\end{minipage}
\item ข\ 
\begin{minipage}[t]{0.9\linewidth}
\raggedright เพชรดูดกลืน UV แล้วเปล่งแสงฟ้า (fluorescence)
\end{minipage}
\item ค\ 
\begin{minipage}[t]{0.9\linewidth}
\raggedright เพชรสะท้อนแสงฟ้าจากทุกทิศทาง
\end{minipage}
\item ง\ 
\begin{minipage}[t]{0.9\linewidth}
\raggedright เพชรเปล่งรังสี X-ray
\end{minipage}
\end{enumerate}
\vspace{-0.3\baselineskip}
\begin{small}
\textit{คำตอบ: ข}
\end{small}

\item ความเร็วของอิเล็กตรอนในวงโคจร n=1 ของอะตอมไฮโดรเจนมีค่าประมาณเท่าใด?
\begin{enumerate}[label=\alph*)., itemsep=0.3\baselineskip, topsep=0.3\baselineskip]
\item ก\ 
\begin{minipage}[t]{0.9\linewidth}
\raggedright 2.2 × 10⁶ m/s
\end{minipage}
\item ข\ 
\begin{minipage}[t]{0.9\linewidth}
\raggedright 2.2 × 10⁷ m/s
\end{minipage}
\item ค\ 
\begin{minipage}[t]{0.9\linewidth}
\raggedright 2.2 × 10⁸ m/s
\end{minipage}
\item ง\ 
\begin{minipage}[t]{0.9\linewidth}
\raggedright 2.2 × 10⁵ m/s
\end{minipage}
\end{enumerate}
\vspace{-0.3\baselineskip}
\begin{small}
\textit{คำตอบ: ก}
\end{small}

\item สเปกตรัมเส้นในช่วง Visible ของไฮโดรเจนมีกี่เส้นเฉพาะที่อยู่ในช่วงที่ตามองเห็นได้?
\begin{enumerate}[label=\alph*)., itemsep=0.3\baselineskip, topsep=0.3\baselineskip]
\item ก\ 
\begin{minipage}[t]{0.9\linewidth}
\raggedright 2 เส้น
\end{minipage}
\item ข\ 
\begin{minipage}[t]{0.9\linewidth}
\raggedright 3 เส้น
\end{minipage}
\item ค\ 
\begin{minipage}[t]{0.9\linewidth}
\raggedright 4 เส้น
\end{minipage}
\item ง\ 
\begin{minipage}[t]{0.9\linewidth}
\raggedright 5 เส้น
\end{minipage}
\end{enumerate}
\vspace{-0.3\baselineskip}
\begin{small}
\textit{คำตอบ: ค}
\end{small}

\item การทดลองใดพิสูจน์ว่าอะตอมไม่ได้เป็นทรงกลมตัน?
\begin{enumerate}[label=\alph*)., itemsep=0.3\baselineskip, topsep=0.3\baselineskip]
\item ก\ 
\begin{minipage}[t]{0.9\linewidth}
\raggedright การเลี้ยวเบนของอิเล็กตรอน
\end{minipage}
\item ข\ 
\begin{minipage}[t]{0.9\linewidth}
\raggedright การกระเจิงของอนุภาคแอลฟา
\end{minipage}
\item ค\ 
\begin{minipage}[t]{0.9\linewidth}
\raggedright การแผ่รังสีของวัตถุดำ
\end{minipage}
\item ง\ 
\begin{minipage}[t]{0.9\linewidth}
\raggedright ปรากฏการณ์โฟโตอิเล็กทริก
\end{minipage}
\end{enumerate}
\vspace{-0.3\baselineskip}
\begin{small}
\textit{คำตอบ: ข}
\end{small}

\item เพราะเหตุใดไฮโดรเจนในธรรมชาติจึงไม่ค่อยเปล่งแสงปกติ แต่เมื่ออยู่ในหลอดไฟนีออนจึงเปล่งแสง?
\begin{enumerate}[label=\alph*)., itemsep=0.3\baselineskip, topsep=0.3\baselineskip]
\item ก\ 
\begin{minipage}[t]{0.9\linewidth}
\raggedright เพราะไฮโดรเจนเป็นก๊าซที่ไม่เปล่งแสง
\end{minipage}
\item ข\ 
\begin{minipage}[t]{0.9\linewidth}
\raggedright เพราะต้องกระตุ้นด้วยไฟฟ้าเพื่อให้อิเล็กตรอนกระโดด
\end{minipage}
\item ค\ 
\begin{minipage}[t]{0.9\linewidth}
\raggedright เพราะไฮโดรเจนต้องอยู่ในสุญญากาศ
\end{minipage}
\item ง\ 
\begin{minipage}[t]{0.9\linewidth}
\raggedright เพราะไฮโดรเจนเปล่งแสงที่ตามองเห็นไม่ได้
\end{minipage}
\end{enumerate}
\vspace{-0.3\baselineskip}
\begin{small}
\textit{คำตอบ: ข}
\end{small}

\item รัศมีโบร์ (Bohr radius) ของอะตอมไฮโดรเจนมีค่าเท่าใด?
\begin{enumerate}[label=\alph*)., itemsep=0.3\baselineskip, topsep=0.3\baselineskip]
\item ก\ 
\begin{minipage}[t]{0.9\linewidth}
\raggedright 0.0529 nm
\end{minipage}
\item ข\ 
\begin{minipage}[t]{0.9\linewidth}
\raggedright 0.529 Å ซึ่งเท่ากับ 0.0529 nm
\end{minipage}
\item ค\ 
\begin{minipage}[t]{0.9\linewidth}
\raggedright 5.29 nm
\end{minipage}
\item ง\ 
\begin{minipage}[t]{0.9\linewidth}
\raggedright 0.0529 pm
\end{minipage}
\end{enumerate}
\vspace{-0.3\baselineskip}
\begin{small}
\textit{คำตอบ: ข}
\end{small}

\item ถ้าต้องการให้อะตอมไฮโดรเจนเปล่งเส้นสเปกตรัมที่ตามองเห็นได้ อิเล็กตรอนต้องถูกกระตุ้นจากสถานะพื้นไปที่ระดับใดขึ้นไป?
\begin{enumerate}[label=\alph*)., itemsep=0.3\baselineskip, topsep=0.3\baselineskip]
\item ก\ 
\begin{minipage}[t]{0.9\linewidth}
\raggedright n = 2
\end{minipage}
\item ข\ 
\begin{minipage}[t]{0.9\linewidth}
\raggedright n = 3
\end{minipage}
\item ค\ 
\begin{minipage}[t]{0.9\linewidth}
\raggedright n = 4
\end{minipage}
\item ง\ 
\begin{minipage}[t]{0.9\linewidth}
\raggedright n = 5
\end{minipage}
\end{enumerate}
\vspace{-0.3\baselineskip}
\begin{small}
\textit{คำตอบ: ข}
\end{small}

\item แสงที่เห็นจากหลอดฟลูออเรสเซนต์มีช่วงความยาวคลื่นต่างจากแสงที่กระตุ้น (UV) อย่างไร?
\begin{enumerate}[label=\alph*)., itemsep=0.3\baselineskip, topsep=0.3\baselineskip]
\item ก\ 
\begin{minipage}[t]{0.9\linewidth}
\raggedright ยาวกว่าเสมอ (Stokes shift)
\end{minipage}
\item ข\ 
\begin{minipage}[t]{0.9\linewidth}
\raggedright สั้นกว่าเสมอ
\end{minipage}
\item ค\ 
\begin{minipage}[t]{0.9\linewidth}
\raggedright เท่าเดิมเสมอ
\end{minipage}
\item ง\ 
\begin{minipage}[t]{0.9\linewidth}
\raggedright ขึ้นอยู่กับชนิดของสาร
\end{minipage}
\end{enumerate}
\vspace{-0.3\baselineskip}
\begin{small}
\textit{คำตอบ: ก}
\end{small}

\end{enumerate}
\end{document}